% Adaptador único del capítulo 38 del sucesor 102.
% Sustituye exclusivamente la jerarquía de encabezados. Los cuerpos de los
% dos propietarios científicos y del delta Ley 9 se conservan línea a línea,
% reordenados conforme al DAG causal aprobado. Véase MAPA_SEGMENTOS_CAPITULO_38.tsv.

% HMT_SOURCE_SEGMENT_BEGIN id=B00 source=colaboracion/parte_iii/source/public_final/base_83/c30_body.tex body_lines=1-2
\label{chap:p3-83-30-angulo_contraangulo}

% HMT_SOURCE_SEGMENT_END id=B00
% HMT_SOURCE_SEGMENT_BEGIN id=A00 source=colaboracion/parte_iii/source/public_final/ascensos_20260826/16f_realizacion_analitica_contraangulo.tex body_lines=1-4
\label{p3-20260826:chap:realizacion-analitica-contraangulo}
\index{contraángulo}
\index{microfibra de pi@microfibra de \(\pi\)!carácter analítico}
\index{carácter determinantal}
% HMT_SOURCE_SEGMENT_END id=A00
% HMT_SOURCE_SEGMENT_BEGIN id=D00 source=manuscrito/sucesor_102/deltas_ley9/c38_par_angular_cartas_y_canales.tex body_lines=1
% Distribución causal exacta del delta Ley 9; contenido conservado sin síntesis.
% HMT_SOURCE_SEGMENT_END id=D00

\section{Fase angular intrínseca \texorpdfstring{\(\alpha_{\mathrm{HMT}}\)}{alpha HMT}, elevación nonádica \texorpdfstring{\(10^3\)}{10^3} y precursor \texorpdfstring{\(x_A\)}{x_A}}
\label{sec:s102:c38:1}

\subsection{Construcción del precursor analítico \texorpdfstring{\(x_A\)}{x_A}, de la contracción \texorpdfstring{\(D_A\)}{D_A} y de la corrección regional \texorpdfstring{\(\mathscr C_\pi(\alpha)\)}{C_pi(alpha)}}
\label{sec:s102:c38:1:1}

% HMT_SOURCE_SEGMENT_BEGIN id=B01 source=colaboracion/parte_iii/source/public_final/base_83/c30_body.tex body_lines=4-32
\label{p3-83-c30-s1:chap:realizacion-analitica-contraangulo}
\index{coordenada angular conjugada \(C^*\)}
\index{microfibra de pi@microfibra de \(\pi\)!carácter analítico}
\index{carácter determinantal}

El catálogo finito y una evaluación analítica responden a preguntas de tipos
distintos. El catálogo determina qué regiones existen, cuáles comparten un
retorno y qué información conserva cada una. Una evaluación analítica debe
decidir, además, cómo se transforma la longitud de una palabra en grado, cómo
se prolonga una rama después de su retorno y con qué carácter escalar se lee
esa prolongación. Este capítulo construye esa segunda operación de manera
explícita.

La construcción parte de cuatro datos ya obtenidos en los capítulos
anteriores: la microfibra de cinco regiones de \(\pi\), la longitud seis de
cada palabra regional, el cierre dodecafásico de \(\alpha\) y el transporte
bicapa asociado a la coordenada común
\(A=10^3\alpha\). Ningún valor de la coordenada angular conjugada \(C_*^{\rm HMT}\), de la componente reducida de
acción, del funcional de Barbero--Immirzi o de una constante metrológica se
emplea en la evaluación. El resultado es un carácter regional convergente,
una recurrencia exacta y una fase orientada cuya carta sexagesimal reproduce
la coordenada angular conjugada publicada.

La separación entre la parte finita y la parte analítica es esencial. El
entero \(169\) será un teorema del catálogo. La serie que lo prolonga será el
teorema de un lector analítico cuyas reglas se declaran antes de evaluarlo.
Así, las premisas adicionales no quedan ocultas dentro de una coincidencia
decimal.

% HMT_SOURCE_SEGMENT_END id=B01
\section{Construcción preangular de \texorpdfstring{\(\eta_{\mathrm{ret}}\)}{eta_ret} mediante \texorpdfstring{\(H_5\)}{H5}, \texorpdfstring{\(D_A\)}{D_A} y \texorpdfstring{\(\mathscr C_\pi(\alpha)\)}{C_pi(alpha)}}
\label{sec:s102:c38:2}

\subsection{Derivación independiente de \texorpdfstring{\(\eta_{\mathrm{ret}}=H_5-(90/\pi)\mathscr C_\pi(\alpha)\)}{eta_ret=H5-(90/pi)C_pi(alpha)} mediante el lector analítico regional}
\label{sec:s102:c38:2:1}

% HMT_SOURCE_SEGMENT_BEGIN id=A01 source=colaboracion/parte_iii/source/public_final/ascensos_20260826/16f_realizacion_analitica_contraangulo.tex body_lines=5-38,40-47

El catálogo finito y una evaluación analítica responden a preguntas de tipos
distintos. El catálogo determina qué regiones existen, cuáles comparten un
retorno y qué información conserva cada una. Una evaluación analítica debe
decidir, además, cómo se transforma la longitud de una palabra en grado, cómo
se prolonga una rama después de su retorno y con qué carácter escalar se lee
esa prolongación. Esta sección construye esa segunda operación de manera
explícita.

La construcción parte de la microfibra de cinco regiones de \(\pi\), la
longitud seis de cada palabra regional, el cierre dodecafásico de \(\alpha\)
y la inserción angular declarada
\[
 \iota_\angle(x):=10^3x\,u_\circ,\qquad
 \Adeg:=10^3\alpha .
\]
Aquí \(u_\circ\) es la base sexagesimal cuya vuelta completa tiene
coordenada \(360\).
El factor \(10^3\) pertenece a esa normalización de carta y no se deduce del
cierre aritmético de \(\alpha\). El transporte bicapa se construye después
a partir de \(\Adeg\). La función que evalúa el lector no recibe como
argumento el contraángulo, la coordenada posterior al retorno, el funcional
hexada--octada ni una constante metrológica. El resultado es un carácter
regional convergente, una recurrencia exacta y una fase orientada en la carta
sexagesimal declarada.

La separación entre la parte finita y la parte analítica es esencial. El
entero \(169\) es un teorema del catálogo. La serie que lo prolonga pertenece
a un lector analítico cuyas reglas se declaran antes de evaluarlo.

La fase se denota por \(y_*\) y, cuando se enfatiza su origen regional, por
\(y_{\rm reg}\); ambos símbolos designan el mismo lector construido a
continuación.


La evaluación recibe el catálogo de \(468\) emisiones, el valor exacto de
\(\alpha_{\rm HMT}\), \(\varphi=(1+\sqrt5)/2\), la fórmula fijada de \(H_5\),
la carta angular y las cinco reglas del lector regional. Con esos datos
calcula \(169\), \(D_A\), la recurrencia y \(C^*\). El teorema es exacto en
esta clase declarada de lectores; no usa \(C^*\) como entrada ni afirma que
las reglas de \(H_5\) sean la única clase natural posible.

% HMT_SOURCE_SEGMENT_END id=A01
\subsection{Persistencia del retorno en las \texorpdfstring{\(5\)}{5} regiones de la microfibra}
\label{sec:s102:c38:2:2}

% HMT_SOURCE_SEGMENT_BEGIN id=B02 source=colaboracion/parte_iii/source/public_final/base_83/c30_body.tex body_lines=34-155

Sea
\[
\mathcal F_\pi=\Psi^{-1}(010211)\subset\mathcal U_6
\]
la microfibra relativa al catálogo. Recordemos sus cinco elementos, ahora
separados en bloque inicial y bloque terminal:
\[
\begin{aligned}
 &(501,614,498\mid169,272,272),\\
 &(810,923,870\mid169,272,272),\\
 &(810,983,810\mid169,272,272),\\
 &(870,923,810\mid169,272,272),\\
 &(870,923,810\mid418,674,521).
\end{aligned}
\]
Definimos la proyección de retorno
\[
p_{\rm ret}:\mathbb Z^6\longrightarrow\mathbb Z^3,
\qquad
p_{\rm ret}(u_1,\ldots,u_6)=(u_4,u_5,u_6).
\]

\begin{proposition}[Bloque terminal persistente]
\label{p3-83-c30-s1:prop:bloque-terminal-persistente}
El multiconjunto \(p_{\rm ret}(\mathcal F_\pi)\) posee un único elemento de
multiplicidad máxima:
\[
b_\pi=(169,272,272),
\qquad
\operatorname{mult}_{\mathcal F_\pi}(b_\pi)=4.
\]
El único bloque terminal restante es
\[
b_\pi^-=(418,674,521)
\]
y tiene multiplicidad uno.
\end{proposition}

\begin{proof}
La afirmación se obtiene por enumeración exacta de las cinco filas del
catálogo que satisfacen \(\Psi(U)=010211\). Cuatro filas poseen bloque
terminal \((169,272,272)\); la quinta, que porta el retorno mutado, posee
\((418,674,521)\). La multiplicidad máxima es cuatro y se alcanza una sola
vez.
\end{proof}

La diferencia orientada entre el retorno mutado y el persistente es
\[
 \omega_\pi=b_\pi^- - b_\pi=(249,402,249),
 \qquad
 \langle(1,1,1),\omega_\pi\rangle=900.
\]
Esta identidad separa dos objetos que no deben confundirse. El entero
\(169\) es una coordenada de un dato terminal, es decir, un dato de grado
cero. El vector \(\omega_\pi\) puede leerse como una cocadena de grado uno
después de orientar la arista que une ambos bloques. Llamarlo cociclo exigiría,
además, fijar un complejo de cocadenas y verificar que su coborde se anula.
La prolongación analítica construida más adelante no presupone esa
identificación.

La selección de \(169\) no exige privilegiar la primera posición del bloque.
El grupo simétrico \(\mathfrak S_3\) actúa permutando sus tres coordenadas.
El estabilizador de \(b_\pi\) intercambia las dos apariciones de \(272\) y
fija la única coordenada de multiplicidad uno. Denotemos por
\(\operatorname{sing}(b)\) esa coordenada cuando el bloque tiene tipo
\((r,s,s)\), con \(r\ne s\). Equivalentemente,
\[
\operatorname{sing}(b)
=\sum_{j=1}^3b_j-2\operatorname{moda}(b).
\]
En particular,
\[
\boxed{\operatorname{sing}(b_\pi)=169.}
\]

\begin{definition}[Selector residual de la microfibra]
\label{p3-83-c30-s1:def:selector-residual-pi}
El selector \(\rho_\pi\) aplica sucesivamente dos operaciones:
\begin{enumerate}
  \item selecciona el bloque terminal de multiplicidad máxima dentro de
  \(p_{\rm ret}(\mathcal F_\pi)\);
  \item retiene la coordenada fija por el estabilizador de ese bloque.
\end{enumerate}
Por la \cref{p3-83-c30-s1:prop:bloque-terminal-persistente},
\[
\boxed{\rho_\pi=169.}
\]
\end{definition}

\begin{theorem}[Unicidad equivariante del retorno]
\label{p3-83-c30-s1:thm:unicidad-retorno-169}
Entre los selectores que
\begin{enumerate}
  \item son invariantes bajo permutaciones de las cinco regiones;
  \item son equivariantes bajo permutaciones de las tres coordenadas
  terminales;
  \item seleccionan el bloque de multiplicidad máxima;
  \item retienen una coordenada fija por su estabilizador,
\end{enumerate}
el valor de retorno está determinado de manera única y es \(169\).
\end{theorem}

\begin{proof}
La invariancia bajo \(\mathfrak S_5\) impide usar el orden de las filas. La
tercera condición selecciona \(b_\pi\), que es el único modo por la
\cref{p3-83-c30-s1:prop:bloque-terminal-persistente}. Su estabilizador en
\(\mathfrak S_3\) posee dos órbitas sobre las coordenadas: la pareja de
valores \(272\) y el valor unipuntual \(169\). La cuarta condición selecciona la
segunda órbita. La equivarianza conserva su valor al desplazar su posición.
\end{proof}

Este argumento mejora la lectura posicional del retorno. El \(169\) no es
«la primera coordenada» de una fila escogida: es un invariante del
multiconjunto regional y de la acción de sus simetrías de reordenación.
Tampoco se lo selecciona por rareza global. En las \(468\) emisiones del
catálogo, \(169\) aparece \(84\) veces, uniformemente distribuido con
\(14\) apariciones en cada una de las seis posiciones. La unicidad del
\cref{p3-83-c30-s1:thm:unicidad-retorno-169} es, por tanto, relacional y regional: reside
en la combinación de microfibra, persistencia modal y estabilizador, no en la
frecuencia aislada del entero.

% HMT_SOURCE_SEGMENT_END id=B02
\subsection{Construcción del lector analítico regional y de su dominio de convergencia}
\label{sec:s102:c38:2:3}

% HMT_SOURCE_SEGMENT_BEGIN id=B05 source=colaboracion/parte_iii/source/public_final/base_83/c30_body.tex body_lines=261-301

La extensión analítica se fija mediante las reglas siguientes.

\begin{definition}[Lector regional determinantal]
\label{p3-83-c30-s1:def:lector-regional-determinantal}
El lector regional determinantal satisface:
\begin{enumerate}
  \item \emph{compatibilidad de grado}: la longitud cilíndrica es el grado
  analítico;
  \item \emph{descomposición de renovación}: el impulso finito de retorno y
  la continuación estricta pertenecen a sumandos aditivos distintos;
  \item \emph{normalización de continuación}: después de registrar una sola
  vez el residuo \(\rho_\pi\), la rama superviviente comienza en la unidad de
  la línea determinante;
  \item \emph{covariancia por concatenación}: cada prolongación elemental
  actúa mediante \(\bigwedge^2\mathcal T=D_A\);
  \item \emph{orientación}: en la convención cardinal fijada para APP, el
  retorno regional pertenece al sector compensador y se sustrae de la fase
  de acción de quinto grado.
\end{enumerate}
\end{definition}

La tercera regla tiene contenido matemático. El valor \(169\) se registra
una vez como amplitud del retorno; no se multiplica de nuevo en cada
prolongación. La continuación estricta cuenta la única línea que prosigue y
la normaliza por su unidad. Esta es la diferencia entre una recurrencia de
renovación y la iteración homogénea del impulso inicial.

La necesidad de declarar esa normalización puede probarse. Para todo
\(\lambda\in\mathbb R\), la familia
\[
F_\lambda(z)
=169z^6+\lambda\frac{D_Az^7}{1-D_Az}
\]
conserva el mismo retorno finito y la misma razón determinantal entre
coeficientes consecutivos. El catálogo y la razón por sí solos no distinguen
\(\lambda\). La unidad de la línea determinante fija
\(\lambda=1\) antes de evaluar \(z=\alpha\). De este modo, la normalización
no se obtiene ajustando el valor final; forma parte de la definición del
carácter.

% HMT_SOURCE_SEGMENT_END id=B05
\subsection{Resolución analítica de la recurrencia de retorno}
\label{sec:s102:c38:2:4}

% HMT_SOURCE_SEGMENT_BEGIN id=B06 source=colaboracion/parte_iii/source/public_final/base_83/c30_body.tex body_lines=303-364

Definimos los coeficientes \(\kappa_n\) por
\[
\kappa_6=169,
\qquad
\kappa_7=D_A,
\qquad
\kappa_{n+1}=D_A\kappa_n\quad(n\ge7).
\]
La primera igualdad es el impulso regional; las dos siguientes describen la
continuación normalizada.

\begin{theorem}[Realización analítica graduada]
\label{p3-83-c30-s1:thm:realizacion-analitica-graduada}
El germen analítico compatible con las reglas de la
\cref{p3-83-c30-s1:def:lector-regional-determinantal} es único y vale
\[
\boxed{
\mathscr C_\pi(z)
=169z^6+\frac{D_Az^7}{1-D_Az}.}
\]
Sus coeficientes satisfacen
\[
\kappa_n=D_A^{n-6}\qquad(n\ge7).
\]
\end{theorem}

\begin{proof}
Escribamos
\[
F(z)=169z^6+\sum_{k\ge1}c_kz^{6+k}.
\]
La normalización de continuación y la covariancia determinantal dan
\[
c_1=D_A,
\qquad
c_{k+1}=D_Ac_k.
\]
Por inducción, \(c_k=D_A^k\). La suma geométrica produce
\[
\sum_{k\ge1}D_A^kz^{6+k}
=\frac{D_Az^7}{1-D_Az}.
\]
No queda ningún coeficiente libre.
\end{proof}

Como \(0<\alpha<1\) y \(0<D_A<1\),
\[
0<D_A\alpha<1.
\]
En consecuencia, \(\mathscr C_\pi(\alpha)\) converge absolutamente y su
radio de convergencia es \(D_A^{-1}>1\). Numéricamente,
\[
D_A=0.7751291192471564301888840964802\ldots,
\]
\[
D_A\alpha=0.00565639046986492673885079095469\ldots.
\]
La rapidez de esta contracción explica por qué la suma completa y su
truncamiento de grado siete poseen las mismas quince primeras cifras en la
carta final.

% HMT_SOURCE_SEGMENT_END id=B06
\subsection{Persistencia del retorno en las \texorpdfstring{\(5\)}{5} regiones y compatibilidad de la componente \texorpdfstring{\(\eta_{\mathrm{ret}}\)}{eta_ret}}
\label{sec:s102:c38:2:5}

% HMT_SOURCE_SEGMENT_BEGIN id=A02 source=colaboracion/parte_iii/source/public_final/ascensos_20260826/16f_realizacion_analitica_contraangulo.tex body_lines=49-170

Sea
\[
\mathcal F_\pi=\Psi^{-1}(010211)\subset\mathcal U_6
\]
la microfibra relativa al catálogo. Recordemos sus cinco elementos, ahora
separados en bloque inicial y bloque terminal:
\[
\begin{aligned}
 &(501,614,498\mid169,272,272),\\
 &(810,923,870\mid169,272,272),\\
 &(810,983,810\mid169,272,272),\\
 &(870,923,810\mid169,272,272),\\
 &(870,923,810\mid418,674,521).
\end{aligned}
\]
Definimos la proyección de retorno
\[
p_{\rm ret}:\mathbb Z^6\longrightarrow\mathbb Z^3,
\qquad
p_{\rm ret}(u_1,\ldots,u_6)=(u_4,u_5,u_6).
\]

\begin{proposition}[Bloque terminal persistente]
\label{p3-20260826:prop:bloque-terminal-persistente}
El multiconjunto \(p_{\rm ret}(\mathcal F_\pi)\) posee un único elemento de
multiplicidad máxima:
\[
b_\pi=(169,272,272),
\qquad
\operatorname{mult}_{\mathcal F_\pi}(b_\pi)=4.
\]
El único bloque terminal restante es
\[
b_\pi^-=(418,674,521)
\]
y tiene multiplicidad uno.
\end{proposition}

\begin{proof}
La afirmación se obtiene por enumeración exacta de las cinco filas del
catálogo que satisfacen \(\Psi(U)=010211\). Cuatro filas poseen bloque
terminal \((169,272,272)\); la quinta, que porta el retorno mutado, posee
\((418,674,521)\). La multiplicidad máxima es cuatro y se alcanza una sola
vez.
\end{proof}

La diferencia orientada entre el retorno mutado y el persistente es
\[
 \omega_\pi=b_\pi^- - b_\pi=(249,402,249),
 \qquad
 \langle(1,1,1),\omega_\pi\rangle=900.
\]
Esta identidad separa dos objetos que no deben confundirse. El entero
\(169\) es una coordenada de un dato terminal, es decir, un dato de grado
cero. El vector \(\omega_\pi\) puede leerse como una cocadena de grado uno
después de orientar la arista que une ambos bloques. Llamarlo cociclo exigiría,
además, fijar un complejo de cocadenas y verificar que su coborde se anula.
La prolongación analítica construida más adelante no presupone esa
identificación.

La selección de \(169\) no exige privilegiar la primera posición del bloque.
El grupo simétrico \(\mathfrak S_3\) actúa permutando sus tres coordenadas.
El estabilizador de \(b_\pi\) intercambia las dos apariciones de \(272\) y
fija la única coordenada de multiplicidad uno. Denotemos por
\(\operatorname{sing}(b)\) esa coordenada cuando el bloque tiene tipo
\((r,s,s)\), con \(r\ne s\). Equivalentemente,
\[
\operatorname{sing}(b)
=\sum_{j=1}^3b_j-2\operatorname{moda}(b).
\]
En particular,
\[
\boxed{\operatorname{sing}(b_\pi)=169.}
\]

\begin{definition}[Selector residual de la microfibra]
\label{p3-20260826:def:selector-residual-pi}
El selector \(\rho_\pi\) aplica sucesivamente dos operaciones:
\begin{enumerate}
  \item selecciona el bloque terminal de multiplicidad máxima dentro de
  \(p_{\rm ret}(\mathcal F_\pi)\);
  \item retiene la coordenada fija por el estabilizador de ese bloque.
\end{enumerate}
Por la \cref{p3-20260826:prop:bloque-terminal-persistente},
\[
\boxed{\rho_\pi=169.}
\]
\end{definition}

\begin{theorem}[Unicidad equivariante del retorno]
\label{p3-20260826:thm:unicidad-retorno-169}
Entre los selectores que
\begin{enumerate}
  \item son invariantes bajo permutaciones de las cinco regiones;
  \item son equivariantes bajo permutaciones de las tres coordenadas
  terminales;
  \item seleccionan el bloque de multiplicidad máxima;
  \item retienen una coordenada fija por su estabilizador,
\end{enumerate}
el valor de retorno está determinado de manera única y es \(169\).
\end{theorem}

\begin{proof}
La invariancia bajo \(\mathfrak S_5\) impide usar el orden de las filas. La
tercera condición selecciona \(b_\pi\), que es el único modo por la
\cref{p3-20260826:prop:bloque-terminal-persistente}. Su estabilizador en
\(\mathfrak S_3\) posee dos órbitas sobre las coordenadas: la pareja de
valores \(272\) y el valor unipuntual \(169\). La cuarta condición selecciona la
segunda órbita. La equivarianza conserva su valor al desplazar su posición.
\end{proof}

Este argumento mejora la lectura posicional del retorno. El \(169\) no es
«la primera coordenada» de una fila escogida: es un invariante del
multiconjunto regional y de la acción de sus simetrías de reordenación.
Tampoco se lo selecciona por rareza global. En las \(468\) emisiones del
catálogo, \(169\) aparece \(84\) veces, uniformemente distribuido con
\(14\) apariciones en cada una de las seis posiciones. La unicidad del
\cref{p3-20260826:thm:unicidad-retorno-169} es, por tanto, relacional y regional: reside
en la combinación de microfibra, persistencia modal y estabilizador, no en la
frecuencia aislada del entero.

% HMT_SOURCE_SEGMENT_END id=A02
\subsection{Aplicación del lector regional al par angular conjugado}
\label{sec:s102:c38:2:6}

% HMT_SOURCE_SEGMENT_BEGIN id=A05 source=colaboracion/parte_iii/source/public_final/ascensos_20260826/16f_realizacion_analitica_contraangulo.tex body_lines=277-317

La extensión analítica se fija mediante las reglas siguientes.

\begin{definition}[Lector regional determinantal]
\label{p3-20260826:def:lector-regional-determinantal}
El lector regional determinantal satisface:
\begin{enumerate}
  \item \emph{compatibilidad de grado}: la longitud cilíndrica es el grado
  analítico;
  \item \emph{descomposición de renovación}: el impulso finito de retorno y
  la continuación estricta pertenecen a sumandos aditivos distintos;
  \item \emph{normalización de continuación}: después de registrar una sola
  vez el residuo \(\rho_\pi\), la rama superviviente comienza en la unidad de
  la línea determinante;
  \item \emph{covariancia por concatenación}: cada prolongación elemental
  actúa mediante \(\bigwedge^2\mathcal T=D_A\);
  \item \emph{orientación}: en la convención cardinal fijada para APP, el
  retorno regional pertenece al sector compensador y se sustrae de la fase
  de acción de quinto grado.
\end{enumerate}
\end{definition}

La tercera regla tiene contenido matemático. El valor \(169\) se registra
una vez como amplitud del retorno; no se multiplica de nuevo en cada
prolongación. La continuación estricta cuenta la única línea que prosigue y
la normaliza por su unidad. Esta es la diferencia entre una recurrencia de
renovación y la iteración homogénea del impulso inicial.

La necesidad de declarar esa normalización puede probarse. Para todo
\(\lambda\in\mathbb R\), la familia
\[
F_\lambda(z)
=169z^6+\lambda\frac{D_Az^7}{1-D_Az}
\]
conserva el mismo retorno finito y la misma razón determinantal entre
coeficientes consecutivos. El catálogo y la razón por sí solos no distinguen
\(\lambda\). La unidad de la línea determinante fija
\(\lambda=1\) antes de evaluar \(z=\alpha\). De este modo, la normalización
no se obtiene ajustando el valor final; forma parte de la definición del
carácter.

% HMT_SOURCE_SEGMENT_END id=A05
\subsection{Suma cerrada y publicación de la componente posterior al retorno}
\label{sec:s102:c38:2:7}

% HMT_SOURCE_SEGMENT_BEGIN id=A06 source=colaboracion/parte_iii/source/public_final/ascensos_20260826/16f_realizacion_analitica_contraangulo.tex body_lines=319-380

Definimos los coeficientes \(\kappa_n\) por
\[
\kappa_6=169,
\qquad
\kappa_7=D_A,
\qquad
\kappa_{n+1}=D_A\kappa_n\quad(n\ge7).
\]
La primera igualdad es el impulso regional; las dos siguientes describen la
continuación normalizada.

\begin{theorem}[Realización analítica graduada]
\label{p3-20260826:thm:realizacion-analitica-graduada}
El germen analítico compatible con las reglas de la
\cref{p3-20260826:def:lector-regional-determinantal} es único y vale
\[
\boxed{
\mathscr C_\pi(z)
=169z^6+\frac{D_Az^7}{1-D_Az}.}
\]
Sus coeficientes satisfacen
\[
\kappa_n=D_A^{n-6}\qquad(n\ge7).
\]
\end{theorem}

\begin{proof}
Escribamos
\[
F(z)=169z^6+\sum_{k\ge1}c_kz^{6+k}.
\]
La normalización de continuación y la covariancia determinantal dan
\[
c_1=D_A,
\qquad
c_{k+1}=D_Ac_k.
\]
Por inducción, \(c_k=D_A^k\). La suma geométrica produce
\[
\sum_{k\ge1}D_A^kz^{6+k}
=\frac{D_Az^7}{1-D_Az}.
\]
No queda ningún coeficiente libre.
\end{proof}

Como \(0<\alpha<1\) y \(0<D_A<1\),
\[
0<D_A\alpha<1.
\]
En consecuencia, \(\mathscr C_\pi(\alpha)\) converge absolutamente y su
radio de convergencia es \(D_A^{-1}>1\). Numéricamente,
\[
D_A=0.7751291192471564301888840964802\ldots,
\]
\[
D_A\alpha=0.00565639046986492673885079095469\ldots.
\]
La rapidez de esta contracción explica por qué la suma completa y su
truncamiento de grado siete poseen las mismas quince primeras cifras en la
carta final.

% HMT_SOURCE_SEGMENT_END id=A06
\section{Confluencia de \texorpdfstring{\(\eta_{\mathrm{ret}}\)}{eta_ret} con la valoración \texorpdfstring{\(\varepsilon_H=-34\)}{epsilon_H=-34} en la magnitud \texorpdfstring{\(\hbar_{\mathrm{ret}}\)}{hbar_ret}}
\label{sec:s102:c38:3}

\subsection{Derivación del carácter determinantal del transporte parangular}
\label{sec:s102:c38:3:1}

% HMT_SOURCE_SEGMENT_BEGIN id=B04 source=colaboracion/parte_iii/source/public_final/base_83/c30_body.tex body_lines=196-259

El cierre dodecafásico proporciona \(\alpha\). La completación de las tres
parejas nonádicas proporciona la escala \(10^3\), de modo que
\[
A=10^3\alpha.
\]
Elegida la carta angular arquimediana, la coordenada común en radianes es
\[
x=\frac{\pi A}{180}=\frac{50\pi}{9}\alpha.
\]
Sea \(R\) una involución real de traza nula en dimensión dos,
\[
R^2=I_2,
\qquad
\operatorname{tr}R=0,
\]
y consideremos el transporte formal bicapa
\[
\mathcal T(x,y)=\exp(-xI_2-yR).
\]
La variable orientada \(y\) todavía no se ha evaluado. Sin embargo, la acción
de \(\mathcal T\) sobre la línea determinante es independiente de ella:
\[
\begin{aligned}
\det\mathcal T(x,y)
&=\exp\!\left(\operatorname{tr}(-xI_2-yR)\right)\\
&=e^{-2x}.
\end{aligned}
\]
Definimos
\[
\boxed{
D_A=e^{-2x}=e^{-100\pi\alpha/9}.}
\]
El subíndice recuerda que se trata del determinante de la contracción común
asociada a \(A\), no de una función de \(C_*^{\rm HMT}\).

\begin{proposition}[Carácter de la línea determinante]
\label{p3-83-c30-s1:prop:caracter-determinante}
La aplicación
\[
\delta_A:(\mathbb N_0,+)\longrightarrow(\mathbb R_{>0},\cdot),
\qquad
\delta_A(k)=D_A^k,
\]
es el único carácter multiplicativo cuyo valor sobre una prolongación
elemental es \(D_A\).
\end{proposition}

\begin{proof}
La multiplicatividad da
\[
\delta_A(k)
=\delta_A(\underbrace{1+\cdots+1}_{k\text{ veces}})
=\delta_A(1)^k=D_A^k.
\]
La misma identidad prueba existencia y unicidad.
\end{proof}

El determinante permanece invariante bajo conjugación de
\(\mathcal T\) y bajo el intercambio de hojas \(R\mapsto-R\). Por tanto, la
cola construida a partir de \(D_A\) pertenece al modo común. La orientación
aparecerá en el signo con el que el carácter regional corrige la fase.

% HMT_SOURCE_SEGMENT_END id=B04
\subsection{Carácter de acción de grado \texorpdfstring{\(5\)}{5} en la construcción regional}
\label{sec:s102:c38:3:2}

% HMT_SOURCE_SEGMENT_BEGIN id=B07 source=colaboracion/parte_iii/source/public_final/base_83/c30_body.tex body_lines=366-405

El bloque central
\[
\mathcal B_c=\begin{pmatrix}7&2\\2&7\end{pmatrix}
\]
posee autovalores \(9\) y \(5\). El ciclo dodecafásico aporta el rango
\(12\); sus órbitas de paso tres poseen longitud cuatro; y el censo completo
aporta el cociente exacto
\[
\frac{104976}{1944}=54.
\]
A partir de estos invariantes se fija el carácter de acción de quinto grado
\[
\boxed{
H_5(\alpha,\varphi)
=\frac12\sqrt{\frac{1000\alpha}{\varphi}}-\alpha
+\frac9{16}\alpha^2-\frac59\alpha^3
+\frac7{48}\alpha^4-\frac1{54}\alpha^5.}
\]
Los cuatro coeficientes racionales pueden escribirse como
\[
\frac9{16}=\frac{\lambda_+}{4^2},
\qquad
\frac59=\frac{\lambda_-}{\lambda_+},
\qquad
\frac7{48}=\frac{\tfrac12\operatorname{tr}\mathcal B_c}{4\cdot12},
\qquad
\frac1{54}=\frac{1944}{104976},
\]
con \((\lambda_+,\lambda_-)=(9,5)\). Estas identidades certifican la
procedencia de los números utilizados. La asignación sucesiva de esos
invariantes a los grados dos, tres, cuatro y cinco forma parte del carácter
analítico; no se deduce de la mera existencia del catálogo finito.

Esta precisión elimina una ambigüedad habitual. Una combinación de
invariantes de la construcción es una definición interna y reproducible; su
naturalidad requiere las reglas del lector que especifican orden, signos y
normalización. Aquí esas reglas son visibles y quedan sometidas a los mismos
controles negativos que el resto de la construcción.

% HMT_SOURCE_SEGMENT_END id=B07
\subsection{Descomposición exacta de la recta de acción en secciones pre-retorno y post-retorno}
\label{sec:s102:c38:3:3}

% HMT_SOURCE_SEGMENT_BEGIN id=B12 source=colaboracion/parte_iii/source/public_final/base_83/c30_body.tex body_lines=571-617

El carácter de quinto grado y la componente de acción posterior al retorno
son dos valores distintos:
\[
H_5
=1.05457181764615446962582182943306\ldots,
\]
\[
\bar h_{\rm HMT}^{\rm ret}
=1.05457181691503906113369058472430\ldots.
\]
El primero es la acción local antes de incorporar el carácter regional; el
segundo es la acción orientada que queda después del retorno. La corrección
que genera la coordenada angular conjugada impide identificarlos.

La comparación metrológica debe respetar esta distinción. El primer valor da
la mantisa
\[
2\pi H_5
=6.62607014999998792600111034989947\ldots,
\]
mientras que la acción posterior da
\[
2\pi\bar h_{\rm HMT}^{\rm ret}
=6.62607014540625433351074984759288\ldots.
\]
En el Sistema Internacional, el valor
\[
h=6.62607015\times10^{-34}\;\mathrm{J\,s}
\]
es exacto por definición \cite{BIPMSI}. Por ello,
\[
2\pi H_5-6.62607015
=-1.2073998889650101\ldots\times10^{-14},
\]
y no es una igualdad exacta. La proximidad del valor anterior al retorno es
un control numérico notable; no autoriza a sustituir la constante definitoria
del SI por la componente posterior ni a confundir las dos acciones HMT\@.

Este resultado determina con precisión la situación matemática. El lector
regional genera \(C_*^{\rm HMT}\) y, desde él, el funcional
\(\Gamma_{6\to8}\). El mismo cálculo separa la aproximación de Planck
producida por \(H_5\) de la acción reducida que interviene en el
la coordenada angular conjugada. La separación no debilita la construcción: elimina una
identificación incompatible con las cifras exactas y convierte cada
comparación en una afirmación falsable de tipo bien definido.

% HMT_SOURCE_SEGMENT_END id=B12
\subsection{Factorización del carácter determinantal en las \texorpdfstring{\(2\)}{2} cartas de fase}
\label{sec:s102:c38:3:4}

% HMT_SOURCE_SEGMENT_BEGIN id=A04 source=colaboracion/parte_iii/source/public_final/ascensos_20260826/16f_realizacion_analitica_contraangulo.tex body_lines=211-275

El cierre dodecafásico proporciona \(\alpha\). La inserción angular declarada
emplea la escala \(10^3\), compatible con la agrupación de tres parejas
nonádicas, de modo que
\[
\Adeg=10^3\alpha\,{}^\circ.
\]
Elegida la carta angular arquimediana, la coordenada común en radianes es
\[
x=: \Arad=\frac{\pi\Adeg}{180}=\frac{50\pi}{9}\alpha.
\]
Sea \(R\) una involución real de traza nula en dimensión dos,
\[
R^2=I_2,
\qquad
\operatorname{tr}R=0,
\]
y consideremos el transporte formal bicapa
\[
\mathcal T(x,y)=\exp(-xI_2-yR).
\]
La variable orientada \(y\) todavía no se ha evaluado. Sin embargo, la acción
de \(\mathcal T\) sobre la línea determinante es independiente de ella:
\[
\begin{aligned}
\det\mathcal T(x,y)
&=\exp\!\left(\operatorname{tr}(-xI_2-yR)\right)\\
&=e^{-2x}.
\end{aligned}
\]
Definimos
\[
\boxed{
D_A=e^{-2x}=e^{-100\pi\alpha/9}.}
\]
El subíndice recuerda que se trata del determinante de la contracción común
asociada a \(A\), no de una función del contraángulo.

\begin{proposition}[Carácter de la línea determinante]
\label{p3-20260826:prop:caracter-determinante}
La aplicación
\[
\delta_A:(\mathbb N_0,+)\longrightarrow(\mathbb R_{>0},\cdot),
\qquad
\delta_A(k)=D_A^k,
\]
es el único carácter multiplicativo cuyo valor sobre una prolongación
elemental es \(D_A\).
\end{proposition}

\begin{proof}
La multiplicatividad da
\[
\delta_A(k)
=\delta_A(\underbrace{1+\cdots+1}_{k\text{ veces}})
=\delta_A(1)^k=D_A^k.
\]
La misma identidad prueba existencia y unicidad.
\end{proof}

El determinante permanece invariante bajo conjugación de
\(\mathcal T\) y bajo el intercambio de hojas \(R\mapsto-R\). Por tanto, la
cola construida a partir de \(D_A\) pertenece al modo común. La orientación
aparecerá en el signo con el que el carácter regional corrige la fase.

% HMT_SOURCE_SEGMENT_END id=A04
\subsection{Publicación del carácter de acción de grado \texorpdfstring{\(5\)}{5} en la carta parangular}
\label{sec:s102:c38:3:5}

% HMT_SOURCE_SEGMENT_BEGIN id=A07 source=colaboracion/parte_iii/source/public_final/ascensos_20260826/16f_realizacion_analitica_contraangulo.tex body_lines=382-418

El bloque central
\[
\mathcal B_c=\begin{pmatrix}7&2\\2&7\end{pmatrix}
\]
posee autovalores \(9\) y \(5\). El ciclo dodecafásico aporta el rango
\(12\); sus órbitas de paso tres poseen longitud cuatro; y el censo completo
aporta el cociente exacto
\[
\frac{104976}{1944}=54.
\]
El lector vigente declara el carácter de acción de quinto grado
\[
\boxed{
H_5(\alpha,\varphi)
=\frac12\sqrt{\frac{1000\alpha}{\varphi}}-\alpha
+\frac9{16}\alpha^2-\frac59\alpha^3
+\frac7{48}\alpha^4-\frac1{54}\alpha^5.}
\]
Los cuatro coeficientes racionales pueden escribirse como
\[
\frac9{16}=\frac{\lambda_+}{4^2},
\qquad
\frac59=\frac{\lambda_-}{\lambda_+},
\qquad
\frac7{48}=\frac{\tfrac12\operatorname{tr}\mathcal B_c}{4\cdot12},
\qquad
\frac1{54}=\frac{1944}{104976},
\]
con \((\lambda_+,\lambda_-)=(9,5)\). Estas identidades fijan la
representabilidad aritmética interna de los coeficientes del carácter. Su
orden, sus signos y su normalización pertenecen a la regla graduada del lector
regional y se aplican antes de toda carta de unidades. Las ablaciones
certifican que cada término interviene efectivamente en la fase resultante.
En particular, ni la mantisa SI de la acción ni el valor sexagesimal del
contraángulo forman parte del dominio de la función generadora.

% HMT_SOURCE_SEGMENT_END id=A07
\section{Aplicación parangular \texorpdfstring{\(P(\alpha,\eta_{\mathrm{ret}})\)}{P(alpha,eta_ret)} y clausura bicapa \texorpdfstring{\(C^*=2(\eta_{\mathrm{ret}}+\alpha)\)}{C*=2(eta_ret+alpha)}}
\label{sec:s102:c38:4}

\subsection{Fase orientada y coordenada angular conjugada}
\label{sec:s102:c38:4:1}

% HMT_SOURCE_SEGMENT_BEGIN id=B08 source=colaboracion/parte_iii/source/public_final/base_83/c30_body.tex body_lines=407-472

La fase de quinto grado anterior al retorno es
\[
y_5=\frac\pi{90}\bigl(H_5+\alpha\bigr).
\]
El carácter regional completo produce la corrección
\(\mathscr C_\pi(\alpha)\). Con la orientación fijada en la
\cref{p3-83-c30-s1:def:lector-regional-determinantal}, definimos
\[
\boxed{
y_*=\frac\pi{90}\bigl(H_5+\alpha\bigr)
-169\alpha^6
-\frac{D_A\alpha^7}{1-D_A\alpha}.}
\]
Ésta es la fórmula primaria. La unidad sexagesimal no interviene hasta aplicar la carta
\[
C_*^{\rm HMT}=\frac{180}{\pi}y_*.
\]

\begin{theorem}[Coordenada angular conjugada del lector regional determinantal]
\label{p3-83-c30-s1:thm:contraangulo-regional}
Fijados el cierre dodecafásico de \(\alpha\),
\(\varphi=(1+\sqrt5)/2\) y las reglas del lector regional determinantal, la
fase \(y_*\) y su coordenada angular conjugada quedan determinados sin introducir el valor
buscado. Su evaluación es
\[
\boxed{
y_*=0.03706622646583826398493779409230638\ldots,}
\]
\[
\boxed{
C_*^{\rm HMT}
=2.12373833896864572426195138039336\ldots{}^\circ.}
\]
Por tanto, a quince cifras decimales,
\[
\boxed{C_*^{\rm HMT}=2.123738338968646^\circ.}
\]
\end{theorem}

\begin{proof}
El \cref{p3-83-c30-s1:thm:unicidad-retorno-169} determina \(169\). El cierre de
\(\alpha\) determina \(A\), \(x\) y
\(D_A=e^{-100\pi\alpha/9}\). El \cref{p3-83-c30-s1:thm:realizacion-analitica-graduada}
determina la serie completa. La fórmula de \(H_5\) contiene únicamente
\(\alpha\), \(\varphi\) e invariantes estructurales declarados. Sustituir
estos datos en la expresión de \(y_*\) da el valor impreso. El valor decimal
de \(C_*^{\rm HMT}\) no aparece en ninguno de los miembros derechos.
\end{proof}

La componente reducida de acción posterior al retorno queda ahora como
salida:
\[
\boxed{
\bar h_{\rm HMT}^{\rm ret}
=\frac{C_*^{\rm HMT}}2-\alpha
=H_5-\frac{90}{\pi}\mathscr C_\pi(\alpha).}
\]
Numéricamente,
\[
\bar h_{\rm HMT}^{\rm ret}
=1.05457181691503906113369058472430\ldots.
\]
La identidad inversa no se ha usado para introducir \(C_*^{\rm HMT}\); se
obtiene después de generar \(y_*\).

% HMT_SOURCE_SEGMENT_END id=B08
\subsection{Coordenada conjugada \texorpdfstring{\(C^*\)}{C*} y reversión de su orientación bajo la involución bicapa}
\label{sec:s102:c38:4:2}

% HMT_SOURCE_SEGMENT_BEGIN id=A08 source=colaboracion/parte_iii/source/public_final/ascensos_20260826/16f_realizacion_analitica_contraangulo.tex body_lines=420-503

La fase de quinto grado anterior al retorno es
\[
y_5=\frac\pi{90}\bigl(H_5+\alpha\bigr).
\]
El carácter regional completo produce la corrección
\(\mathscr C_\pi(\alpha)\). Con la orientación fijada en la
\cref{p3-20260826:def:lector-regional-determinantal}, definimos
\[
\boxed{
\Crad:=y_*=\frac\pi{90}\bigl(H_5+\alpha\bigr)
-169\alpha^6
-\frac{D_A\alpha^7}{1-D_A\alpha}.}
\]
En adelante escribimos también
\[
 \boxed{y_{\rm reg}:=y_*}
\]
cuando sea necesario distinguir esta realización de la fase espectral de
\Gtr. Las dos notaciones designan aquí el mismo objeto regional canónico.
Ésta es la fórmula primaria. La unidad sexagesimal no interviene hasta aplicar la carta
\[
\boxed{\Cdeg:=\frac{180}{\pi}\Crad.}
\]
Los símbolos \(\Crad\) y \(\Cdeg\) no designan dos contraángulos: son las
coordenadas radial y sexagesimal del mismo elemento angular.

\begin{theorem}[Contraángulo del lector regional determinantal]
\label{p3-20260826:thm:contraangulo-regional}
Fijados el cierre dodecafásico de \(\alpha\),
\(\varphi=(1+\sqrt5)/2\), el germen \(H_5\) y las reglas del lector regional
determinantal, la fase \(\Crad\) y su contraángulo quedan determinados sin
introducir el valor buscado en la función de evaluación. Su evaluación es
\[
\boxed{
\Crad=0.03706622646583826398493779409230638\ldots,}
\]
\[
\boxed{
\Cdeg
=2.12373833896864572426195138039336\ldots{}^\circ.}
\]
Por tanto, a quince cifras decimales,
\[
\boxed{\Cdeg=2.123738338968646^\circ.}
\]
\end{theorem}

\begin{proof}
El \cref{p3-20260826:thm:unicidad-retorno-169} determina \(169\). El cierre de
\(\alpha\) determina \(\Adeg\), \(\Arad\) y
\(D_A=e^{-100\pi\alpha/9}\). El \cref{p3-20260826:thm:realizacion-analitica-graduada}
determina la serie completa. La fórmula fijada de \(H_5\) contiene únicamente
\(\alpha\), \(\varphi\) e invariantes estructurales declarados. Sustituir
estos datos en la expresión de \(\Crad\) da el valor impreso. El valor decimal
del contraángulo no aparece en ninguno de los miembros derechos de esta
evaluación. La comparación con una mantisa metrológica de acción se efectúa
únicamente después de obtener la fase y no modifica la función generadora.
\end{proof}

Sea \(u_\circ\) la base sexagesimal cuya vuelta completa tiene coordenada
\(360\), e introduzcamos la coordenada angular
\(\alpha_{\angle,{\rm deg}}:=\alpha_{\rm HMT}\). La salida posterior al
retorno queda entonces tipada por
\[
\boxed{
\etaRet
=\frac{\Cdeg}2-\alpha_{\angle,{\rm deg}}
=H_5-\frac{90}{\pi}\mathscr C_\pi(\alpha).}
\]
Su coordenada radial es
\[
 \etaRetrad=\frac{\pi}{180}\etaRet
 =\frac{\Crad}{2}-\frac{\pi}{180}\alpha_{\rm HMT}.
\]
Numéricamente,
\[
\etaRet
=1.05457181691503906113369058472430\ldots.
\]
La identidad inversa no se ha usado para introducir \(\Cdeg\); se obtiene
después de generar \(\Crad\). El símbolo \(\hbar\) queda reservado para el
elemento de la recta de acción construido en la sección siguiente.

% HMT_SOURCE_SEGMENT_END id=A08
\subsection{Teorema directo \texorpdfstring{\(C^*=2(\eta_{\mathrm{ret}}+\alpha)\)}{C*=2(eta_ret+alpha)} y carácter inverso de la identidad \texorpdfstring{\(\eta_{\mathrm{ret}}=C^*/2-\alpha\)}{eta_ret=C*/2-alpha}}
\label{sec:s102:c38:4:3}

% HMT_SOURCE_SEGMENT_BEGIN id=A13 source=colaboracion/parte_iii/source/public_final/ascensos_20260826/16f_realizacion_analitica_contraangulo.tex body_lines=750-797

El carácter de quinto grado y la coordenada posterior al retorno
son dos valores distintos:
\[
H_5
=1.05457181764615446962582182943306\ldots,
\]
\[
\etaRet
=1.05457181691503906113369058472430\ldots.
\]
El primero es el carácter local antes de incorporar el carácter regional; el
segundo es la coordenada orientada que queda después del retorno. La corrección
que genera el contraángulo impide identificarlos.

La comparación metrológica debe respetar esta distinción. El primer valor da
la mantisa
\[
2\pi H_5
=6.62607014999998792600111034989947\ldots,
\]
mientras que la coordenada posterior da
\[
2\pi\etaRet
=6.62607014540625433351074984759288\ldots.
\]
Como control externo únicamente, retenemos la mantisa \(6.62607015\) del
valor definitorio del Sistema Internacional \cite{BIPMSI}. No se incorpora
todavía su exponente ni su unidad a la recta HMT: la década se deriva en el
\cref{chap:orden-determinantal-accion} y la unidad pertenece a una carta
metrológica posterior. La comparación de mantisas da
\[
2\pi H_5-6.62607015
=-1.2073998889650101\ldots\times10^{-14},
\]
y no es una igualdad exacta. La proximidad del valor anterior al retorno es
un control numérico notable; no autoriza a sustituir la constante definitoria
del SI por la coordenada posterior ni a confundir una coordenada angular con
el elemento de acción que se construirá al incorporar la década.

Este resultado determina con precisión la situación matemática. El lector
regional genera \((\Crad,\Cdeg)\) y, desde ese único elemento angular, el funcional
\(\Gamma_{6\to8}\). El mismo cálculo separa la proximidad externa de la
mantisa producida por \(H_5\) de la coordenada de retorno que interviene en el
contraángulo. La separación no debilita la construcción: elimina una
identificación incompatible con las cifras exactas y convierte cada
comparación en una afirmación falsable de tipo bien definido.

% HMT_SOURCE_SEGMENT_END id=A13
% HMT_SOURCE_SEGMENT_BEGIN id=D01 source=manuscrito/sucesor_102/deltas_ley9/c38_par_angular_cartas_y_canales.tex body_lines=3-20
\label{sec:l9s102:par-angular}

La completación cubica
\begin{equation}
 10^3=(9+1)^3=9^3+3\cdot9^2+3\cdot9+1
 =729+243+27+1
 \label{eq:l9s102:base-1000}
\end{equation}
determina la carta global de la coordenada directa. La clausura bicapa aporta
el factor \(2\) de la coordenada conjugada:
\begin{equation}
 \boxed{
 A_{\deg}=10^3\alpha_{\mathrm{HMT}},
 \qquad
 C^*_{\deg}=2(\eta_{\mathrm{ret}}+\alpha_{\mathrm{HMT}})
 \quad\text{en }\mathbb R/(360\mathbb Z).}
 \label{eq:l9s102:a-cstar-deg}
\end{equation}
% HMT_SOURCE_SEGMENT_END id=D01
% HMT_SOURCE_SEGMENT_BEGIN id=D04 source=manuscrito/sucesor_102/deltas_ley9/c38_par_angular_cartas_y_canales.tex body_lines=60-66
La identidad inversa
\begin{equation}
 \eta_{\mathrm{ret}}=\frac12C^*_{\deg}-\alpha_{\mathrm{HMT}}
 \label{eq:l9s102:eta-corolario-inverso}
\end{equation}
queda establecida como corolario algebraico de la construcción directa.

% HMT_SOURCE_SEGMENT_END id=D04
\section{Conmutatividad exacta de las cartas en vueltas, grados y radianes}
\label{sec:s102:c38:5}

\subsection{Derivación de la longitud cilíndrica y del grado analítico del precursor parangular}
\label{sec:s102:c38:5:1}

% HMT_SOURCE_SEGMENT_BEGIN id=B03 source=colaboracion/parte_iii/source/public_final/base_83/c30_body.tex body_lines=157-194

Sea \(\mathcal P\) el módulo libre generado por las palabras regionales y
sea \(F^n\mathcal P\) su filtración por longitud. La graduación asociada
registra el primer nivel en el que aparece cada palabra. Para un parámetro
\(z\), el carácter ordinario de longitud es
\[
\chi_z([w])=z^{|w|}.
\]
La multiplicatividad
\[
\chi_z([uv])=\chi_z([u])\chi_z([v])
\]
expresa que la concatenación suma longitudes. Al evaluar en
\(z=\alpha\), una palabra de longitud \(n\) recibe grado \(\alpha^n\).

\begin{definition}[Compatibilidad de grado]
\label{p3-83-c30-s1:def:compatibilidad-grado}
El lector analítico regional es compatible con la filtración cilíndrica
cuando envía el componente homogéneo de longitud \(n\) al grado analítico
\(n\):
\[
\operatorname{gr}_n\mathcal P\longrightarrow
\alpha^n\mathbb R.
\]
\end{definition}

Las cinco regiones pertenecen a \(\mathcal U_6\); por ello, el dato de
retorno aparece en grado seis:
\[
\rho_\pi[U_6]\longmapsto169\alpha^6.
\]
Este seis es la longitud de la palabra regional. No es la longitud de los
paseos cerrados del \cref{chap:cinco-ciclos-pi}. Aquellos paseos recorren a
la vez los cuatro anclajes
\(U_\pi,U_e,U_\varphi,U_{\rm APP}\) y poseen longitud mínima ocho. Una
palabra, una arista del grafo de bloques y un paseo cerrado son objetos de
tipos distintos; la graduación utilizada aquí pertenece al primero.

% HMT_SOURCE_SEGMENT_END id=B03
\subsection{Publicación del grado analítico en la carta global de fase}
\label{sec:s102:c38:5:2}

% HMT_SOURCE_SEGMENT_BEGIN id=A03 source=colaboracion/parte_iii/source/public_final/ascensos_20260826/16f_realizacion_analitica_contraangulo.tex body_lines=172-209

Sea \(\mathcal P\) el módulo libre generado por las palabras regionales y
sea \(F^n\mathcal P\) su filtración por longitud. La graduación asociada
registra el primer nivel en el que aparece cada palabra. Para un parámetro
\(z\), el carácter ordinario de longitud es
\[
\chi_z([w])=z^{|w|}.
\]
La multiplicatividad
\[
\chi_z([uv])=\chi_z([u])\chi_z([v])
\]
expresa que la concatenación suma longitudes. Al evaluar en
\(z=\alpha\), una palabra de longitud \(n\) recibe grado \(\alpha^n\).

\begin{definition}[Compatibilidad de grado]
\label{p3-20260826:def:compatibilidad-grado}
El lector analítico regional es compatible con la filtración cilíndrica
cuando envía el componente homogéneo de longitud \(n\) al grado analítico
\(n\):
\[
\operatorname{gr}_n\mathcal P\longrightarrow
\alpha^n\mathbb R.
\]
\end{definition}

Las cinco regiones pertenecen a \(\mathcal U_6\); por ello, el dato de
retorno aparece en grado seis:
\[
\rho_\pi[U_6]\longmapsto169\alpha^6.
\]
Este seis es la longitud de la palabra regional. No es la longitud de los
paseos cerrados del \cref{chap:cinco-ciclos-pi}. Aquellos paseos recorren a
la vez los cuatro anclajes
\(U_\pi,U_e,U_\varphi,U_{\rm APP}\) y poseen longitud mínima ocho. Una
palabra, una arista del grafo de bloques y un paseo cerrado son objetos de
tipos distintos; la graduación utilizada aquí pertenece al primero.

% HMT_SOURCE_SEGMENT_END id=A03
% HMT_SOURCE_SEGMENT_BEGIN id=D02 source=manuscrito/sucesor_102/deltas_ley9/c38_par_angular_cartas_y_canales.tex body_lines=21-34
El calendario global exhibe las tres factorizaciones
\begin{equation}
 360=12\cdot30=4\cdot90=3\cdot120.
 \label{eq:l9s102:360-factorizaciones}
\end{equation}
La carta analítica se obtiene mediante
\begin{equation}
 A_{\mathrm{rad}}=\frac\pi{180}A_{\deg}
 =\frac{50\pi}{9}\alpha_{\mathrm{HMT}},
 \qquad
 C^*_{\mathrm{rad}}=\frac\pi{180}C^*_{\deg}
 =\frac\pi{90}(\eta_{\mathrm{ret}}+\alpha_{\mathrm{HMT}}).
 \label{eq:l9s102:a-cstar-rad}
\end{equation}
% HMT_SOURCE_SEGMENT_END id=D02
% HMT_SOURCE_SEGMENT_BEGIN id=D03 source=manuscrito/sucesor_102/deltas_ley9/c38_par_angular_cartas_y_canales.tex body_lines=35-59

\begin{theorem}[Conmutatividad de las factorizaciones regional y de acción]
\label{thm:l9s102:conmutatividad-cstar}
Las dos rutas hacia la fase conjugada producen la misma salida:
\begin{equation}
 \boxed{
 \frac\pi{180}C^*_{\deg}
 =\frac\pi{90}(\eta_{\mathrm{ret}}+\alpha_{\mathrm{HMT}})
 =\frac\pi{90}(H_5+\alpha_{\mathrm{HMT}})
 -\mathscr C_\pi(\alpha_{\mathrm{HMT}})=y_*.}
 \label{eq:l9s102:cuadrado-cstar}
\end{equation}
\end{theorem}

\begin{proof}
Sustituyendo \cref{eq:l9s102:eta-ret} en el segundo miembro se obtiene
\[
 \frac\pi{90}
 \left(H_5-\frac{90}{\pi}\mathscr C_\pi+\alpha_{\mathrm{HMT}}\right)
 =\frac\pi{90}(H_5+\alpha_{\mathrm{HMT}})-\mathscr C_\pi.
\]
La última expresion es la definicion regional de \(y_*\), mientras la
primera igualdad procede de \cref{eq:l9s102:a-cstar-deg}.
\end{proof}

% HMT_SOURCE_SEGMENT_END id=D03
\section{Descomposición espectral \texorpdfstring{\(q_\pm\)}{q+/-} y cálculo funcional sobre los canales \texorpdfstring{\(90/120\)}{90/120} previamente generados}
\label{sec:s102:c38:6}

\subsection{Reconstrucción espectral del par angular y transición hexada--octada}
\label{sec:s102:c38:6:1}

% HMT_SOURCE_SEGMENT_BEGIN id=B11 source=colaboracion/parte_iii/source/public_final/base_83/c30_body.tex body_lines=523-569

La fase común y la fase orientada determinan los dos canales
\[
q_-=e^{-x+y_*},
\qquad
q_+=e^{-x-y_*}.
\]
Se recuperan exactamente los invariantes
\[
q_-q_+=D_A,
\qquad
\alpha=-\frac9{100\pi}\log D_A,
\qquad
C_*^{\rm HMT}=\frac{90}{\pi}\log\frac{q_-}{q_+}.
\]
La primera igualdad lee la contracción común; la tercera, la separación
orientada de las hojas. El intercambio \(R\mapsto-R\) permuta
\(q_+\) y \(q_-\), conserva \(D_A\) e invierte el signo de \(y_*\).

Para
\[
S_{90}(q)=\frac{q^{90}}{1-q^{270}},
\qquad
S_{120}(q)=\frac{q^{120}}{1-q^{360}},
\]
el funcional de transición hexada--octada se escribe íntegramente en la
carta de fases:
\[
\Gamma_{6\to8}
=\frac{180}{\pi}xy_*
+12\bigl(S_{90}(q_-)-S_{90}(q_+)\bigr)
-\bigl(S_{120}(q_-)-S_{120}(q_+)\bigr).
\]
La evaluación producida por la coordenada \(C_*^{\rm HMT}\) del
\cref{p3-83-c30-s1:thm:contraangulo-regional} es
\[
\boxed{
\Gamma_{6\to8}
=0.274007672694459274747255260774\ldots.}
\]
Así, el valor interno reconocido en Mecánica Dimensional como funcional de
Barbero--Immirzi deja de recibir \(C_*^{\rm HMT}\) como literal independiente:
ambos descienden del mismo par de fases generado en este capítulo. La
interpretación geométrica y física de \(\Gamma_{6\to8}\) pertenece a la obra
de Mecánica Dimensional; su evaluación matemática pertenece ya a la interfaz
previamente construida.

% HMT_SOURCE_SEGMENT_END id=B11
\subsection{Transportador relativo y disjunción espectral de los bloques}
\label{sec:s102:c38:6:2}

% HMT_SOURCE_SEGMENT_BEGIN id=B13 source=colaboracion/parte_iii/source/public_final/base_83/c30_body.tex body_lines=619-675
\label{p3-83-c30-s1:sec:transportador-relativo-split}

En la base ordenada de canales, sea
\[
 R=\begin{pmatrix}0&1\\1&0\end{pmatrix},
 \qquad P_\pm=\frac{I\pm R}{2},
 \qquad
 B_0=7I+2R,
 \qquad B_1=\Adeg I+\Cdeg R.
\]
La descomposición escindida
\[
 xI+yR=\rho(x,y)e^{\eta(x,y)R},
 \quad \rho(x,y)=\sqrt{x^2-y^2},
 \quad \eta(x,y)=\operatorname{artanh}(y/x)
\]
construye el transportador único
\[
 T_{\rm rel}
 =\frac{\sqrt{(\Adeg)^2-(\Cdeg)^2}}{\sqrt{45}}
   \exp\!\left[
   \left(\operatorname{artanh}\frac{\Cdeg}{\Adeg}
   -\operatorname{artanh}\frac27\right)R\right],
 \qquad T_{\rm rel}B_0=B_1.
\]
En la base espectral actúa por
\(9\mapsto\Adeg+\Cdeg\) y
\(5\mapsto\Adeg-\Cdeg\).

\begin{theorem}[Disjunción espectral de las realizaciones local y angular]
\label{p3-83-c30-s1:thm:no-go-entrelazador-bloques-split}
Los espectros
\[
 \operatorname{Spec}(B_0)=\{9,5\},\qquad
 \operatorname{Spec}(B_1)=\{\Adeg+\Cdeg,\Adeg-\Cdeg\}
\]
son disjuntos. En consecuencia,
\[
 \operatorname{Hom}_{B_0,B_1}
 :=\{\Phi:\Phi B_0=B_1\Phi\}=\{0\}.
\]
\end{theorem}

\begin{proof}
Las cotas racionales construidas para las coordenadas angulares dan
\(7.29<\Adeg<7.30\) y \(2.12<\Cdeg<2.13\); por tanto,
\(9.41<\Adeg+\Cdeg<9.43\) y
\(5.16<\Adeg-\Cdeg<5.18\). En la base que diagonaliza \(R\), cada entrada
de un entrelazador queda multiplicada por la diferencia entre un autovalor
de \(B_0\) y uno de \(B_1\); la disjunción obliga a que todas se anulen.
\end{proof}

El transportador relativo compara dos bloques dentro del álgebra escindida;
los entrelazadores dodecafásico--Witt y de incidencia actúan entre
representaciones equivalentes en sus residencias propias. Esta separación
tipa positivamente ambas operaciones y conserva sus dominios.

% HMT_SOURCE_SEGMENT_END id=B13
\subsection{Evaluación espectral \texorpdfstring{\(V_{90,120}\)}{V_90,120} de los canales de incidencia previamente generados por el TPK}
\label{sec:s102:c38:6:3}

% HMT_SOURCE_SEGMENT_BEGIN id=A11 source=colaboracion/parte_iii/source/public_final/ascensos_20260826/16f_realizacion_analitica_contraangulo.tex body_lines=554-572

La fase común y la fase orientada determinan los dos canales
\[
q_-=e^{-\Arad+\Crad},
\qquad
q_+=e^{-\Arad-\Crad}.
\]
Se recuperan exactamente los invariantes
\[
q_-q_+=D_A,
\qquad
\alpha=-\frac9{100\pi}\log D_A,
\qquad
\Cdeg=\frac{90}{\pi}\log\frac{q_-}{q_+}.
\]
La primera igualdad lee la contracción común; la tercera, la separación
orientada de las hojas. El intercambio \(R\mapsto-R\) permuta
\(q_+\) y \(q_-\), conserva \(D_A\) e invierte el signo de \(\Crad\).

% HMT_SOURCE_SEGMENT_END id=A11
% HMT_SOURCE_SEGMENT_BEGIN id=D05 source=manuscrito/sucesor_102/deltas_ley9/c38_par_angular_cartas_y_canales.tex body_lines=68-97
\label{sec:l9s102:qpm-posteriores}

Sea \(H^2=I\) y \(P_\pm=(I\pm H)/2\). El bloque angular
\begin{equation}
 B=A_{\mathrm{rad}}I+C^*_{\mathrm{rad}}H
 =(A_{\mathrm{rad}}+C^*_{\mathrm{rad}})P_+
 +(A_{\mathrm{rad}}-C^*_{\mathrm{rad}})P_-
 \label{eq:l9s102:b-angular}
\end{equation}
se publica multiplicativamente por
\begin{equation}
 e^{-B}=q_+P_++q_-P_-,
 \qquad
 q_\pm=\exp[-(A_{\mathrm{rad}}\pm C^*_{\mathrm{rad}})].
 \label{eq:l9s102:qpm-angular}
\end{equation}
La involución \(C^*\mapsto-C^*\) permuta \(q_+\) y \(q_-\), y conserva
\begin{equation}
 q_+q_-=e^{-2A_{\mathrm{rad}}}=D_A.
 \label{eq:l9s102:det-angular}
\end{equation}
Los cardinales \(90/120\), construidos previamente por la incidencia del
TPK, reciben ahora su evaluación espectral:
\begin{equation}
 V_{90,120}
 =(I-T^{90})(I-T^{120})^{-1}
 =\sum_{\sigma\in\{+,-\}}
 \frac{1-q_\sigma^{90}}{1-q_\sigma^{120}}P_\sigma.
 \label{eq:l9s102:v90120-posterior}
\end{equation}
% HMT_SOURCE_SEGMENT_END id=D05
\subsubsection{Transportador relativo escindido}
\label{sec:s102:c38:6:3:1}

% HMT_SOURCE_SEGMENT_BEGIN id=A12 source=colaboracion/parte_iii/source/public_final/ascensos_20260826/16f_realizacion_analitica_contraangulo.tex body_lines=574-748
\label{p3-20260826:sec:transportador-relativo-split}

La comparación con el bloque central de APP puede efectuarse sin identificar
representaciones distintas. Fijemos la base ordenada de canales y la
involución
\[
 R=\begin{pmatrix}0&1\\1&0\end{pmatrix},
 \qquad R^2=I,
 \qquad
 P_\pm=\frac{I\pm R}{2}.
\]
En el álgebra conmutativa
\[
 \mathscr B_R:=\{xI+yR:x,y\in\mathbb R\}
\]
consideremos, dentro de la carta sexagesimal ya declarada,
\[
 B_0:=7I+2R,
 \qquad
 B_1:=\Adeg I+\Cdeg R.
\]
El cambio ortogonal de base
\[
 S=\frac1{\sqrt2}
 \begin{pmatrix}1&1\\1&-1\end{pmatrix}
\]
diagonaliza simultáneamente la involución y ambos bloques. Para \(x>|y|\),
definamos
\[
 \rho(x,y):=\sqrt{x^2-y^2},
 \qquad
 \eta(x,y):=\operatorname{artanh}\!\left(\frac yx\right).
\]
Entonces
\[
 xI+yR=\rho(x,y)e^{\eta(x,y)R}.
\]
Si
\[
 \rho_0=\sqrt{45},\quad
 \eta_0=\operatorname{artanh}(2/7),\quad
 \rho_1=\rho(\Adeg,\Cdeg),\quad
 \eta_1=\eta(\Adeg,\Cdeg),
\]
la única multiplicación izquierda que lleva \(B_0\) a \(B_1\) es
\[
 \boxed{
 T_{\rm rel}
 :=\frac{\rho_1}{\sqrt{45}}
 e^{(\eta_1-\eta_0)R},
 \qquad
 T_{\rm rel}B_0=B_1.}
\]

\begin{proposition}[Transportador relativo de los dos bloques]
\label{p3-20260826:prop:transportador-relativo-split}
Una vez fijados \(R\), el orden de los canales, la rama positiva y la carta
sexagesimal, \(T_{\rm rel}\) es el único elemento de
\(\mathscr B_R\) que satisface \(TB_0=B_1\). En la base espectral actúa por
\[
 9\longmapsto\Adeg+\Cdeg,
 \qquad
 5\longmapsto\Adeg-\Cdeg.
\]
\end{proposition}

\begin{proof}
La descomposición polar escindida da la fórmula impresa y prueba la existencia.
Como \(\det B_0=45\ne0\), toda solución matricial de \(TB_0=B_1\) satisface
\(T=B_1B_0^{-1}\); por tanto es única. La diagonalización mediante \(S\)
produce los dos cocientes espectrales
\((\Adeg+\Cdeg)/9\) y \((\Adeg-\Cdeg)/5\), equivalentes a la expresión
exponencial de \(T_{\rm rel}\).
\end{proof}

Este transportador no es un entrelazador representacional. En efecto, un
entrelazador \(\Phi\) entre los endomorfismos \(B_0\) y \(B_1\) debería
satisfacer
\[
 \Phi B_0=B_1\Phi.
\]

\begin{theorem}[Obstrucción espectral al entrelazamiento]
\label{p3-20260826:thm:no-go-entrelazador-bloques-split}
Para los valores de \(\Adeg\) y \(\Cdeg\) construidos en esta sección,
\[
 \boxed{\Phi B_0=B_1\Phi\quad\Longrightarrow\quad\Phi=0.}
\]
\end{theorem}

\begin{proof}
Los espectros son
\[
 \operatorname{Spec}(B_0)=\{9,5\},
 \qquad
 \operatorname{Spec}(B_1)
 =\{\Adeg+\Cdeg,\Adeg-\Cdeg\}.
\]
Las cotas racionales previamente demostradas proporcionan los intervalos
disjuntos
\[
 7.29<\Adeg<7.30,
 \qquad
 2.12<\Cdeg<2.13,
\]
y, por consiguiente,
\[
 9.41<\Adeg+\Cdeg<9.43,
 \qquad
 5.16<\Adeg-\Cdeg<5.18.
\]
En la base que diagonaliza \(R\), cada entrada \(\phi_{ij}\) de \(\Phi\)
queda multiplicada por la diferencia entre un autovalor de \(B_0\) y uno de
\(B_1\). Todas esas diferencias son no nulas; luego todas las entradas de
\(\Phi\) se anulan.
\end{proof}

Numéricamente,
\[
 \rho_1=6.9814819335172378048\ldots,
 \qquad
 \frac{\rho_1}{\sqrt{45}}=1.0407378791354140779\ldots,
\]
\[
 \eta_1-\eta_0
 =0.005796351470571295590\ldots.
\]
La fórmula conserva el álgebra escindida, los proyectores \(P_\pm\) y la forma
lorentziana hasta el factor conforme \(\rho_1^2/45\). Depende, sin embargo,
de \(B_1\) ya construido: no constituye una derivación anterior de
\(\Adeg\) o \(\Cdeg\), ni debe confundirse con los entrelazadores
equivariantes entre realizaciones de un mismo módulo.

Para
\[
S_{90}(q)=\frac{q^{90}}{1-q^{270}},
\qquad
S_{120}(q)=\frac{q^{120}}{1-q^{360}},
\]
el funcional de transición hexada--octada se escribe íntegramente en la
carta de fases:
\[
\Gamma_{6\to8}
=\frac{180}{\pi}\Arad\Crad
+12\bigl(S_{90}(q_-)-S_{90}(q_+)\bigr)
-\bigl(S_{120}(q_-)-S_{120}(q_+)\bigr).
\]
La evaluación producida por el contraángulo del
\cref{p3-20260826:thm:contraangulo-regional} es
\[
\boxed{
\Gamma_{6\to8}
=0.274007672694459274747255260774\ldots.}
\]
Las formas equivalentes del funcional matemático interno son
\[
\boxed{
\Gamma_{6\to8}
=\Arad\Cdeg+K_{6\to8}
=\frac{180}{\pi}\Arad\Crad+K_{6\to8}
=\frac{\pi\Adeg\Cdeg}{180}+K_{6\to8}.}
\]
Estas identidades evalúan \(\Gamma_{6\to8}\) después de generar
\((\Crad,\Cdeg)\); no se utiliza en sentido inverso para seleccionar el
contraángulo. La lectura física en la conexión de Ashtekar--Barbero se
representa después mediante
\[
 \Gamma_{6\to8}\dashrightarrow\gamma_{\rm BI}.
\]
La flecha no introduce un tercer funcional ni retroactúa sobre \(C^*\). Las
series \(S_{90}\), \(S_{120}\), los coeficientes \(12\) y \(-1\) y su
asignación al paso \(6\to8\) pertenecen a la construcción matemática
declarada. La identificación física exige, además, la normalización propia
de la conexión.

% HMT_SOURCE_SEGMENT_END id=A12
\section{Naturalidad bajo refinamiento y separación respecto del reconocimiento metrológico}
\label{sec:s102:c38:7}

\subsection{Estimación exacta del resto de la prolongación analítica}
\label{sec:s102:c38:7:1}

% HMT_SOURCE_SEGMENT_BEGIN id=B09 source=colaboracion/parte_iii/source/public_final/base_83/c30_body.tex body_lines=474-499

La expresión hallada inicialmente,
\[
y_7=\frac\pi{90}(H_5+\alpha)-169\alpha^6-D_A\alpha^7,
\]
es el truncamiento de grado siete del carácter completo. Produce
\[
C_7=2.12373833896864600265403827103919\ldots{}^\circ.
\]
El resto omitido no es un término indeterminado; es la cola geométrica
exacta
\[
\mathscr C_\pi(\alpha)
-169\alpha^6-D_A\alpha^7
=\frac{D_A^2\alpha^8}{1-D_A\alpha}.
\]
Por tanto,
\[
\left|C_7-C_*^{\rm HMT}\right|
=\frac{180}{\pi}\frac{D_A^2\alpha^8}{1-D_A\alpha}
=2.7839208689064583\ldots\times10^{-16}\;{}^\circ.
\]
La recurrencia transforma así una aproximación de séptimo grado en una
realización analítica completa y proporciona una cota cerrada para cada
truncamiento posterior.

% HMT_SOURCE_SEGMENT_END id=B09
\subsection{Dependencia funcional del selector regional y de la continuación analítica}
\label{sec:s102:c38:7:2}

% HMT_SOURCE_SEGMENT_BEGIN id=B10 source=colaboracion/parte_iii/source/public_final/base_83/c30_body.tex body_lines=501-521

La construcción cambia al retirar cualquiera de sus componentes activos.
Los siguientes valores se obtienen sin recalibrar los demás términos:
\[
\begin{array}{@{}lc@{}}
\toprule
\text{variante}&\text{coordenada angular conjugada en grados}\\
\midrule
\rho_\pi=168&2.1237383389772976863904487\ldots\\
\rho_\pi=169&2.1237383389686457242619514\ldots\\
\rho_\pi=170&2.1237383389599937621334541\ldots\\
\text{sin cola determinantal}&2.1237383389686949415301675\ldots\\
D_A\text{ sustituido por }1&2.1237383389686313409965826\ldots\\
\bottomrule
\end{array}
\]
Estas ablaciones no constituyen por sí solas una estimación probabilística.
Demuestran que el selector equivariante y el carácter determinantal son
componentes efectivas de la salida y que no pueden suprimirse manteniendo la
misma realización.

% HMT_SOURCE_SEGMENT_END id=B10
\subsection{Interpretación operatoria de la coordenada angular conjugada}
\label{sec:s102:c38:7:3}

% HMT_SOURCE_SEGMENT_BEGIN id=B14 source=colaboracion/parte_iii/source/public_final/base_83/c30_body.tex body_lines=677-715

El procedimiento completo puede resumirse como una cadena de aplicaciones:
\[
\begin{aligned}
\mathcal F_\pi
&\longmapsto b_\pi
\longmapsto\rho_\pi=169,\\
\alpha
&\longmapsto A
\longmapsto x
\longmapsto D_A,\\
(\rho_\pi,D_A)
&\longmapsto\mathscr C_\pi(\alpha),\\
(H_5,\mathscr C_\pi(\alpha))
&\longmapsto y_*
\longmapsto C_*^{\rm HMT},\\
(x,y_*)
&\longmapsto(q_+,q_-)
\longmapsto\Gamma_{6\to8}.
\end{aligned}
\]
La primera línea pertenece al catálogo finito; la segunda, al cierre común de
\(\alpha\); la tercera, a la prolongación determinantal; la cuarta produce la
coordenada orientada; la quinta la transporta hacia la interfaz dimensional.

En esta organización, \(\alpha\) mide la contracción común del transporte y
\(C_*^{\rm HMT}\) mide la separación orientada de sus dos hojas. El retorno
regional no añade una constante externa: selecciona una amplitud finita y la
prolonga mediante el determinante ya fijado por \(\alpha\). La carta en
grados es posterior. El objeto primario es el par de fases \((x,y_*)\).

Una verificación exhaustiva reconstruye las cinco regiones directamente
desde el catálogo, recorre las acciones de
\(\mathfrak S_5\) y \(\mathfrak S_3\), comprueba la recurrencia y suma la
serie. La entrada de \(\alpha\) es la salida exacta de su cierre
dodecafásico, no una cifra metrológica transcrita. La comprobación mantiene
separados los teoremas del catálogo, las reglas del lector y el contraste
externo; por ello ninguna comparación experimental interviene en la función
generadora.
% HMT_SOURCE_SEGMENT_END id=B14
\subsection{Certificación de convergencia del truncamiento parangular}
\label{sec:s102:c38:7:4}

% HMT_SOURCE_SEGMENT_BEGIN id=A09 source=colaboracion/parte_iii/source/public_final/ascensos_20260826/16f_realizacion_analitica_contraangulo.tex body_lines=505-530

La expresión hallada inicialmente,
\[
y_7=\frac\pi{90}(H_5+\alpha)-169\alpha^6-D_A\alpha^7,
\]
es el truncamiento de grado siete del carácter completo. Produce
\[
C_7=2.12373833896864600265403827103919\ldots{}^\circ.
\]
El resto omitido no es un término indeterminado; es la cola geométrica
exacta
\[
\mathscr C_\pi(\alpha)
-169\alpha^6-D_A\alpha^7
=\frac{D_A^2\alpha^8}{1-D_A\alpha}.
\]
Por tanto,
\[
\left|C_7-\Cdeg\right|
=\frac{180}{\pi}\frac{D_A^2\alpha^8}{1-D_A\alpha}
=2.7839208689064583\ldots\times10^{-16}\;{}^\circ.
\]
La recurrencia transforma así una aproximación de séptimo grado en una
realización analítica completa y proporciona una cota cerrada para cada
truncamiento posterior.

% HMT_SOURCE_SEGMENT_END id=A09
\subsection{Naturalidad del selector regional bajo refinamiento y cambio de carta}
\label{sec:s102:c38:7:5}

% HMT_SOURCE_SEGMENT_BEGIN id=A10 source=colaboracion/parte_iii/source/public_final/ascensos_20260826/16f_realizacion_analitica_contraangulo.tex body_lines=532-552

La construcción cambia al retirar cualquiera de sus componentes activos.
Los siguientes valores se obtienen sin recalibrar los demás términos:
\[
\begin{array}{@{}lc@{}}
\toprule
\text{variante}&\text{contraángulo en grados}\\
\midrule
\rho_\pi=168&2.1237383389772976863904487\ldots\\
\rho_\pi=169&2.1237383389686457242619514\ldots\\
\rho_\pi=170&2.1237383389599937621334541\ldots\\
\text{sin cola determinantal}&2.1237383389686949415301675\ldots\\
D_A\text{ sustituido por }1&2.1237383389686313409965826\ldots\\
\bottomrule
\end{array}
\]
Estas ablaciones no constituyen por sí solas una estimación probabilística.
Demuestran que el selector equivariante y el carácter determinantal son
componentes efectivas de la salida y que no pueden suprimirse manteniendo la
misma realización.

% HMT_SOURCE_SEGMENT_END id=A10
\subsection{Independencia metrológica de la construcción parangular}
\label{sec:s102:c38:7:6}

% HMT_SOURCE_SEGMENT_BEGIN id=A14 source=colaboracion/parte_iii/source/public_final/ascensos_20260826/16f_realizacion_analitica_contraangulo.tex body_lines=799-831

El procedimiento completo puede resumirse como una cadena de aplicaciones:
\[
\begin{aligned}
\mathcal F_\pi
&\longmapsto b_\pi
\longmapsto\rho_\pi=169,\\
\alpha
&\longmapsto \Adeg
\longmapsto \Arad
\longmapsto D_A,\\
(\rho_\pi,D_A)
&\longmapsto\mathscr C_\pi(\alpha),\\
(H_5,\mathscr C_\pi(\alpha))
&\longmapsto \Crad
\longmapsto \Cdeg,\\
(\Arad,\Crad)
&\longmapsto(q_+,q_-)
\longmapsto\Gamma_{6\to8}.
\end{aligned}
\]
La primera línea pertenece al catálogo finito; la segunda, al cierre común de
\(\alpha\); la tercera, a la prolongación determinantal; la cuarta produce la
coordenada orientada; la quinta la transporta hacia la interfaz dimensional.

En esta organización, \(\alpha\) mide la contracción común del transporte y
el contraángulo mide la separación orientada de sus dos hojas. En la
evaluación vigente, el retorno regional selecciona una amplitud finita y la
prolonga mediante el determinante ya fijado por \(\alpha\), sin insertar una
constante externa adicional. La carta en grados es posterior. El objeto
primario es el par de fases \((\Arad,\Crad)\). La recurrencia, la suma
regional y las ablaciones dependen del catálogo y de
\(\alpha_{\rm HMT}\), pero no reciben una constante metrológica adicional.
% HMT_SOURCE_SEGMENT_END id=A14

